\booksection{נספח לפרק 1: נושא מערכת ההפעלה HarmonyOS}

\chapterintro{HarmonyOS היא מערכת הפעלה בפיתוח עצמי של Huawei, מאז השקתה ב-2019 עברה מספר גרסאות עיקריות. היא אינה רק “עוד תחליף Android”, אלא מין חדש עם תפיסת עיצוב ייחודית. בואו נצלול לעומק מערכת ההפעלה המתפתחת הזו.}

\booksubsection{מהי HarmonyOS? מדוע היא שונה?}

תפיסת הליבה של HarmonyOS היא “חיבור כל הדברים”. בניגוד למערכות הפעלה מסורתיות המשרתות בעיקר מכשיר בודד, HarmonyOS תוכננה מלכתחילה לשיתוף פעולה בין מכשירים.

\textbf{ארכיטקטורה מבוזרת}היא הנשמה של HarmonyOS. בפשטות, היא יכולה להפוך מספר מכשירים ל“מסוף-על” וירטואלי אחד. הטלפון, הטאבלט, המחשב, השעון, הטלוויזיה שלך - אינם עוד ישויות מבודדות, אלא “איברים” של שלם אחד, שיכולים לקרוא זה ליכולותיו של זה.

לדוגמה: אתה צופה בסרטון בטלפון, רוצה לראות על מסך גדול, הדרך המסורתית היא שיקוף; ב-HarmonyOS, אתה יכול ישירות “להזרים” את הסרטון לטלוויזיה ולהמשיך לצפות, הטלפון הופך לשלט. זה לא שיקוף פשוט, אלא היישום עצמו רץ על פני מכשירים.

\booksubsection{HarmonyOS 4.3: גרסה בשלה המבוססת על AOSP}

HarmonyOS 4.3 היא הגרסה הפרוסה באופן נרחב בטלפוני Huawei הנוכחיים. יש לציין ש-HarmonyOS 4.3 תואמת למסגרת הקוד של Android Open Source Project (AOSP), ולכן היא תואמת באופן מלא ליישומי Android - רוב האפליקציות שאתה מוריד מ-Huawei AppGallery הן במהותן עדיין יישומי Android, הרצים דרך שכבת התאימות של HarmonyOS. אך HarmonyOS 4.3 אינה פשוט מעטפת ל-Android, היא מוסיפה על בסיס AOSP כמות גדולה של טכנולוגיות בפיתוח עצמי של Huawei כמו יכולות מבוזרות, מנוע Ark וכו'.

\subsubsection*{תכונות ליבה:}

\begin{itemize}
  \item \textbf{תאימות יישומי Android}: מכיוון שפותחה על בסיס AOSP, HarmonyOS 4.3 יכולה להריץ כמעט את כל יישומי Android, המשתמש כמעט שאינו מרגיש הבדל.
  \item \textbf{מנוע Ark}: מהדר וסביבת ריצה בפיתוח עצמי של Huawei, מאפשרים ליישומים לרוץ חלק יותר ולחסוך בחשמל. בהשוואה לגרסאות מוקדמות, שיפור ניכר בחלקות.
  \item \textbf{מסוף-על}: במרכז הבקרה, ניתן לראות באופן אינטואיטיבי מכשירי HarmonyOS אחרים בסביבה, גרור כדי להתחבר. הטלפון משתמש במקלדת ובעכבר של המחשב, הטאבלט משתמש במצלמת הטלפון - הכל בנוחות.
  \item \textbf{כרטיסי שירות}: דומים לווידג'טים של iOS, אך גמישים יותר. ניתן להחליק למעלה על סמל אפליקציה כדי להרחיב כרטיס המציג מידע חיוני, ללא צורך לפתוח את האפליקציה.
  \item \textbf{מומלץ בחום לכל המשתמשים להפעיל}: מאפשר התקנת יישומים רק מחנות האפליקציות הרשמית, מפחית סיכון לתוכנות זדוניות.
  \item \textbf{מרכז פרטיות}: ניהול מרוכז של הרשאות יישומים, צפייה ברישום התנהגות יישומים.
\end{itemize}

\subsubsection*{מאפייני ממשק משתמש:}

\begin{itemize}
  \item \textbf{מסך נעילה בסגנון מגזין}: מסך הנעילה מציג תמונות ומאמרים יפים, החלק הצידה כדי להיכנס למסך הבית.
  \item \textbf{כרטיסים אוניברסליים}: סמלי מסך הבית ניתנים להרחבה לכרטיסים המציגים מידע נוסף.
  \item \textbf{תחנת מעבר-על}: אזור אחסון זמני, ניתן לאחסן טקסט, תמונות, קבצים באופן זמני, ואז לגרור לכל אפליקציה לשימוש. שיתוף בין יישומים הופך לחלק להפליא.
  \item \textbf{חיפוש חכם}: חיפוש גלובלי, ניתן לחפש יישומים, קבצים, הגדרות, דפי אינטרנט.
\end{itemize}

\booksubsection{HarmonyOS NEXT (HarmonyOS 6): נקודת ההתחלה האמיתית של מערכת ההפעלה העצמאית של Huawei}

HarmonyOS NEXT (HarmonyOS NEXT 5.0) שונה באופן מהותי מ-HarmonyOS 4.3 - היא משתמשת במיקרו-קרנל HarmonyOS בפיתוח עצמי של Huawei, אינה מסתמכת עוד על ליבת Linux וקוד AOSP, וגם\textbf{אינה תומכת עוד ישירות ביישומי Android}. זהו הצעד המכריע של HarmonyOS מ“צורת מעבר תואמת Android” ל“מערכת הפעלה עצמאית לחלוטין”.

\subsubsection*{שינויים מרכזיים:}

\begin{itemize}
  \item \textbf{קרנל עצמאי לחלוטין}: HarmonyOS NEXT משתמשת בקרנל בפיתוח עצמי של Huawei, אינה מכילה עוד קוד AOSP. המערכת רזה יותר, בטוחה יותר, יעילה יותר.
  \item \textbf{אינה תומכת עוד ביישומי Android}: זהו השינוי הגדול ביותר. HarmonyOS NEXT יכולה להריץ רק יישומי HarmonyOS מקוריים, אינה יכולה להתקין ולהריץ קובצי APK. לפני השדרוג אנא ודא שלאפליקציות הנפוצות שלך יש כבר גרסת HarmonyOS מקורית.
  \item \textbf{אקוסיסטם יישומים מקוריים}: Huawei פועלת במלוא המרץ לדחוף מפתחים להעביר יישומים לפלטפורמת HarmonyOS המקורית. נכון לעכשיו, אפליקציות מרכזיות כמו WeChat, Alipay, Taobao וכו' כבר השיקו גרסאות HarmonyOS מקוריות, אך ייתכן שחלק מהיישומים הנישתיים טרם הותאמו.
  \item \textbf{ארכיטקטורת אבטחה Star Shield}: מערכת אבטחה חדשה לחלוטין, הגנה מקצה לקצה מחומרה לתוכנה. תומכת בבקרת הרשאות ברמת דיוק גבוהה יותר, כמו "מיקום מטושטש" - האפליקציה יודעת רק איפה אתה בערך, לא את המיקום המדויק.
  \item \textbf{שילוב עמוק של AI}: יכולות AI ברמת המערכת, לא רק לאפליקציה מסוימת יש AI. למשל כתוביות AI יכולות לתרגם וידאו בזמן אמת, חיתוך AI יכול לזהות נושאים באופן חכם.
  \item \textbf{העברה חלקה בין מכשירים}: יכולות מבוזרות מתקדמות יותר. מצב היישום יכול לעבור באופן חלק בין מכשירים, המסמך שאתה עורך בטלפון נפתח אוטומטית במחשב כשאתה ניגש אליו.
  \item \textbf{שירותי אטום}: שירותים קלי משקל הניתנים לשימוש ללא התקנה. למשל סריקת קוד QR פותחת ישירות פונקציה מסוימת, ללא צורך בהורדת אפליקציה תחילה.
\end{itemize}

\begin{tcbwarn}{חובה לקרוא לפני שדרוג ל-HarmonyOS NEXT}
לאחר שדרוג ל-HarmonyOS NEXT, הטלפון שלך לא יוכל להריץ יישומי Android. אנא ודא לפני השדרוג: (1) לאפליקציות הנפוצות שלך יש כבר גרסת HarmonyOS מקורית; (2) אין בעבודה או בחייך תרחישים המחייבים הסתמכות על יישומי Android. לא ניתן לחזור ל-HarmonyOS 4.3 לאחר השדרוג.

\end{tcbwarn}

\subsubsection*{שינויי לוגיקת תפעול:}

\begin{itemize}
  \item \textbf{ניווט במחוות}: החלק מקצה ימין/שמאל לחזרה, החלק למעלה מלמטה למסך הבית, החלק למעלה ועצור לריבוי משימות.
  \item \textbf{חלונות חכמים}: תומך במסך מפוצל, חלון צף, תצוגה מקבילה ועוד צורות ריבוי משימות.
  \item \textbf{שיחות Changlian}: שירות שיחות הווידאו של Huawei, תומך בשיתוף מסך, יופי, טשטוש רקע.
\end{itemize}

\booksubsection{אקוסיסטם HarmonyOS: מכשירים ויישומים}

HarmonyOS אינה רק מערכת טלפון, אלא אקוסיסטם המכסה מגוון מכשירים.

\subsubsection*{סוגי מכשירים נתמכים:}

\begin{itemize}
  \item \textbf{טלפון וטאבלט}: מכשירי ליבה, הפונקציונליות המלאה ביותר.
  \item \textbf{שעונים חכמים וצמידים}: פועלים באופן עצמאי או בשיתוף עם הטלפון.
  \item \textbf{מסך חכם (טלוויזיה)}: מרכז בידור ביתי.
  \item \textbf{רמקול חכם}: נקודת כניסה לאינטראקציה קולית.
  \item \textbf{מכונית}: תא נהג HarmonyOS, חוויית תא נהג חכם.
  \item \textbf{קישוריות לרכב}: מקרר, מזגן, מכונת כביסה וכו'.
\end{itemize}

\subsubsection*{מצב אקוסיסטם היישומים:}

חנות האפליקציות של HarmonyOS מכילה מיליוני אפליקציות, רוב האפליקציות המרכזיות זמינות. Huawei מקדמת באופן פעיל מפתחים ליצור אפליקציות מקוריות ל-HarmonyOS. נכון לאמצע 2026, WeChat, Alipay, Taobao, Douyin, Meituan, JD ואפליקציות מובילות אחרות כבר השיקו גרסאות מקוריות ל-HarmonyOS, והשימוש היומיומי חלק ברוב המקרים. אפליקציות נישה מסוימות עשויות עדיין להיות בתהליך התאמה, מומלץ לוודא שהאפליקציות השגורות שלך נתמכות לפני שדרוג.

\booksubsection{טיפים מעשיים ל-HarmonyOS}

\begin{itemize}
  \item \textbf{מסוף-על}: החלק מטה מפינה ימנית עליונה של המסך לפתיחת מרכז הבקרה, לחץ על "מסוף-על", ניתן לראות מכשירים אחרים בסביבה, גרור לחיבור.
  \item \textbf{תחנת מעבר-על}: לחץ לחיצה ארוכה על טקסט/תמונה/קובץ, בחר "העתק לתחנת מעבר", ואז באפליקציה אחרת לחץ לחיצה ארוכה להדבקה, או גרור החוצה.
  \item \textbf{כרטיסי שירות}: החלק למעלה על סמל אפליקציה להרחבת כרטיס, לחץ לחיצה ארוכה על הכרטיס לשינוי גודל או הוספה למסך הבית.
  \item \textbf{חלונות חכמים}: החלק פנימה מקצה ימין או שמאל של המסך ועצור, פתח סרגל צד, לחץ על אפליקציה לפתיחה בחלון צף.
  \item \textbf{הגנת פרטיות}: הגדרות > פרטיות, ניתן לצפות ברישום שימוש בהרשאות יישומים, להגדיר מיקום מטושטש וכו'.
\end{itemize}

\begin{tcbwarn}{הערות חשובות}
מערכת HarmonyOS מתעדכנת ומשתדרגת כל הזמן, וייתכנו הבדלים בין גרסאות שונות. תיאור זה מבוסס על המאפיינים העיקריים של HarmonyOS 4.3 ו-6. אנא הסתמך על הגרסה בפועל בטלפון שלך. מומלץ לעדכן את המערכת באופן שוטף כדי לקבל את התכונות העדכניות ועדכוני האבטחה.

\end{tcbwarn}

\begin{tcbpractice}{נסה עכשיו: חווה את התכונות הייחודיות של HarmonyOS}
\begin{enumerate}
  \item פתח את מרכז הבקרה, מצא את "מסוף-על", ובדוק אם יש מכשירי HarmonyOS אחרים בסביבה.
  \item נסה את "תחנת מעבר-על": העתק טקסט לתחנת המעבר, ולאחר מכן גרור אותו החוצה באפליקציית התזכירים.
  \item החלק למעלה על כמה סמלי אפליקציות נפוצות, וראה איזה תוכן מציגים כרטיסי השירות שלהם.
  \item עבור אל הגדרות > פרטיות, ובדוק אילו אפליקציות השתמשו לאחרונה בהרשאות רגישות.
\end{enumerate}

\end{tcbpractice}
